\section{Theory}
\label{sec:theory}

The measurement of \cref{sec:results} reports a dispersion. The theory explains it for the rung whose estimator has a closed form, the single loss approximation of \cref{eq:sla}, and gives two consequences that can be tested: the dispersion predicted at the true parameters of each cell, and the share of histories that conclude the wrong thing about the mean.
Results marked as propositions are derived under the assumptions below; the finite-sample behaviour that the derivation does not cover is stated as a conjecture with its evidence.

\subsection{Assumptions}
\label{sec:theory-assumptions}

\begin{assumption}
\label{as:model}
The recorded losses of a history are $n$ independent draws from a generalised Pareto distribution with tail index $\xi>-\tfrac12$, scale $\beta_H$ and location $H$, and the number of recorded losses is Poisson with mean $k=\lambda_H T$, independent of the amounts.
\end{assumption}

\Cref{as:model} holds for the generator of this work: a generalised Pareto distribution of location $u_0$, truncated from below at $H\ge u_0$, is again generalised Pareto with the same index, location $H$ and scale $\beta_H=\beta+\xi(H-u_0)$, and the recorded rate is $\lambda_H=\lambda(1+\xi(H-u_0)/\beta)^{-1/\xi}$. It does not hold for a real bank, whose severity is not exactly Pareto above the threshold.
The estimator is the maximum likelihood fit $(\hat\xi,\hat\beta_H)$ of the recorded excesses over $H$ with the recorded rate $\hat\lambda_H=N/T$, inserted in \cref{eq:sla}. Write $A=\lambda_H/(1-\alpha)$, $L=\ln A$, $E=A^{\xi}$ and $q=H+(\beta_H/\xi)(E-1)$ for the value of \cref{eq:sla} at the true parameters.

\begin{lemma}[Smith, 1987; Drees, Ferreira and de Haan, 2004]
\label{lem:mle}
Under \cref{as:model}, as $k\to\infty$, $\sqrt k\,\bigl(\hat\xi-\xi,\ \hat\beta_H/\beta_H-1\bigr)$ converges in distribution to a bivariate normal law with mean zero and covariance
\begin{equation}
\Sigma=(1+\xi)\begin{pmatrix}1+\xi&-1\\-1&2\end{pmatrix}.
\label{eq:mle-cov}
\end{equation}
\end{lemma}

The lemma is the standard result for the maximum likelihood estimator of the generalised Pareto parameters, valid for $\xi>-\tfrac12$ and therefore through $\xi=1$ \citep{smith1987,drees2004}; it is a known result and we do not prove it. In particular $\Var(\hat\xi)\approx(1+\xi)^2/k$, so at $\xi=0.96$ with $k=\exK{}$ recorded losses the standard deviation of the fitted index is $(1+\xi)/\sqrt k=\exSdXi{}$.

\subsection{The dispersion of the capital}
\label{sec:theory-dispersion}

\begin{proposition}[Dispersion of the single loss approximation]
\label{prop:dispersion}
Under \cref{as:model}, $\sqrt k\,(\ln\hat q_\alpha-\ln q)$ converges in distribution to a normal law with mean zero and variance
\begin{equation}
k\sigma^2=g_\xi^{2}\,(1+\xi)^2-2\,g_\xi g_\beta\,(1+\xi)+2\,g_\beta^{2}\,(1+\xi)+g_\lambda^{2},
\label{eq:dispersion}
\end{equation}
where the sensitivities of the logarithm of the capital to the index, to the logarithm of the scale and to the logarithm of the rate are
\begin{equation}
g_\xi=\frac{\beta_H}{q}\Bigl[\frac{E\ln A}{\xi}-\frac{E-1}{\xi^2}\Bigr],\qquad g_\beta=\frac{\beta_H(E-1)}{\xi\,q},\qquad g_\lambda=\frac{\beta_H E}{q}.
\label{eq:sensitivities}
\end{equation}
\end{proposition}

\begin{proof}
The capital is the smooth function $q(\xi,\beta_H,\lambda_H)=H+(\beta_H/\xi)\bigl((\lambda_H/(1-\alpha))^{\xi}-1\bigr)$. Differentiating $\ln q$ gives $\partial_\xi\ln q=g_\xi$, $\partial_{\ln\beta}\ln q=g_\beta$ and $\partial_{\ln\lambda}\ln q=g_\lambda$: the last follows from $\partial_\lambda\bigl(A^{\xi}\bigr)=\xi A^{\xi}/\lambda$, which makes $\lambda\,\partial_\lambda q=\beta_H E$. By \cref{lem:mle} and the identity $\ln\hat\beta_H-\ln\beta_H=(\hat\beta_H/\beta_H-1)+o_p(k^{-1/2})$, the pair $(\hat\xi,\ln\hat\beta_H)$ is asymptotically normal with covariance $\Sigma/k$. The count $N$ is Poisson with mean $k$, so $\sqrt k(\ln\hat\lambda_H-\ln\lambda_H)\to\mathcal N(0,1)$, and it is independent of the amounts. The delta method applied to $\ln q$ gives a normal limit with variance $(g_\xi,g_\beta)\Sigma(g_\xi,g_\beta)^{\top}/k+g_\lambda^2/k$, which expands to \cref{eq:dispersion}.
\end{proof}

\begin{corollary}[Leading order]
\label{cor:leading}
If $H\ll q$ and $E\gg1$, as at the level 0.999 with hundreds of recorded losses, then $g_\beta\to1$, $g_\xi\to L-1/\xi$ and $g_\lambda\to\xi$, and
\begin{equation}
k\sigma^2\approx(1+\xi)\Bigl[(1+\xi)\,\ell^{2}-2\,\ell+2\Bigr]+\xi^{2},\qquad \ell=L-\frac1\xi .
\label{eq:leading}
\end{equation}
The dispersion of the logarithm of the capital is of order $(1+\xi)\,\ell/\sqrt k$: it grows with the logarithm of the extrapolation factor $\lambda_H/(1-\alpha)$, that is, with the level, and it falls as $k^{-1/2}$, that is, as the inverse square root of the number of recorded losses.
\end{corollary}

For the generator at $\xi=0.96$, $\lambda_H=73.4$ and $T=5$, one has $L=\ln(\exLam{}/0.001)=\exL{}$, $\ell=\exEll{}$ and $k=\exK{}$, and \cref{eq:leading} gives $\sigma\approx\exSigmaLead{}$, a dispersion of the logarithm of the capital of almost one natural unit: the capital is multiplied or divided by $e^{\sigma}=\exFactor{}$ at one standard deviation. \Cref{tab:theory} evaluates the exact \cref{eq:dispersion} in the 16 cells and compares it with the measurement.

\begin{corollary}[History needed for a stability limit]
\label{cor:history}
If $\ln\hat q-\ln q$ is approximately normal with standard deviation $\sigma$ and the median of the ratio is one, the interquartile range of the ratio is $I=2\sinh(z_{0.75}\sigma)$ with $z_{0.75}=0.6745$. The limit $I^\ast$ is reached when $\sigma^\ast=\operatorname{arsinh}(I^\ast/2)/z_{0.75}$, and since $\sigma^2$ falls as $1/T$,
\begin{equation}
T^\ast=T\,\frac{\sigma^2}{\sigma^{\ast2}},\qquad \sigma^\ast=\frac{\operatorname{arsinh}(I^\ast/2)}{z_{0.75}},
\label{eq:years}
\end{equation}
which for the green limit $I^\ast=0.25$ of the product gives $\sigma^\ast=\sigmaStar{}$.
\end{corollary}

\subsection{Scale invariance and the currency}
\label{sec:theory-invariance}

\begin{proposition}[The dispersion is a property of $(\xi,k,L)$]
\label{prop:invariance}
Let the excesses over $H$ be generalised Pareto with index $\xi$ and scale $\beta_H$. The distribution of $(\hat q-H)/(q-H)$ depends on $\beta_H$ neither through the currency nor through the size of the bank: it depends only on $\xi$, on the number $k$ of recorded losses and on the extrapolation factor $A=\lambda_H/(1-\alpha)$.
\end{proposition}

\begin{proof}
The maximum likelihood estimators are equivariant under a change of unit: if every excess is multiplied by $c>0$, $\hat\xi$ is unchanged and $\hat\beta_H$ is multiplied by $c$. The excess of the capital over $H$ in \cref{eq:sla} is proportional to $\hat\beta_H/\hat\xi$, so it is multiplied by $c$, as is the true excess $q-H$, and the ratio is unchanged. The count enters through $\hat\lambda_H$ only, whose law depends on $k$ and not on the units.
\end{proof}

The proposition says that the currency and the units play no role, and that a bank with twice the rate of losses and the same tail index differs from another only through $k$ and $A$: its dispersion falls as $k^{-1/2}$ while $L$ grows only with the logarithm of the rate, so size helps, but slowly.

\subsection{The wrong conclusion about the mean}
\label{sec:theory-mean}

\begin{proposition}[Share of histories that conclude an infinite mean]
\label{prop:mean}
Under \cref{as:model}, as $k\to\infty$ with $\xi$ fixed,
\begin{equation}
\Prob(\hat\xi\ge1)\to\Phi\!\Bigl(-\frac{(1-\xi)\sqrt k}{1+\xi}\Bigr)\quad\text{for }\xi\neq1,\qquad\Prob(\hat\xi\ge1)\to\tfrac12\ \text{ for }\xi=1 .
\label{eq:misclassification}
\end{equation}
At a finite $k$ the normal approximation $\Phi\bigl(-(1-\xi)\sqrt k/(1+\xi)\bigr)$ is the predicted share, and the indices for which a history cannot decide whether the mean exists, in the sense that the true index lies within $z_{0.975}$ standard deviations of the boundary, are those with $|1-\xi|<z_{0.975}(1+\xi)/\sqrt k$.
\end{proposition}

\begin{proof}
By \cref{lem:mle}, $\sqrt k(\hat\xi-\xi)/(1+\xi)\to\mathcal N(0,1)$. The event $\hat\xi\ge1$ is $\sqrt k(\hat\xi-\xi)/(1+\xi)\ge\sqrt k(1-\xi)/(1+\xi)$, and the limit follows. For $\xi=1$ the right-hand side is zero and the probability tends to one half.
\end{proof}

For $\xi=0.96$ and $k=\exK{}$ the indistinguishability band is $|1-\xi|<\exBandC{}$, so the index $0.96$ and the boundary of one cannot be told apart by the maximum likelihood fit of a five-year history; at $k=\exKEleven{}$ (ten years at $\xi=1.1$, threshold EUR 10,000) the band is $\exBandD{}$. The proposition does not say that the fitted index is unbiased in samples of this size.

The proposition has a practical reading. The share is a decreasing function of $\sqrt k$ for any index below one and an increasing function of it for any index above one, so a longer history sharpens the conclusion about the mean only for a process whose index is clearly on one side of the boundary; for a process near the boundary, the case that matters for a validator because the capital is the most sensitive there, the share stays near one half whatever the history.
The same quantity gives the length of history that would separate an index from the boundary with a stated confidence: it is the $k$ at which the band of indistinguishability shrinks below the distance of the index from one.

\subsection{Conditions under which the results hold}
\label{sec:theory-conditions}

Four conditions bound the propositions. First, $k$ must be large enough for the normal approximation; here $k$ ranges from \kMin{} to \kMax{}, and the approximation of the share is checked, not assumed, in \cref{sec:results}. Second, the family must be correct: the propositions apply to the family of the generator, and outside \cref{as:model} they do not apply; a misspecified fit adds a bias that they do not contain. Third, the single loss approximation targets $q$ only asymptotically in the level:
its value at the true parameters is below the Fourier quantile by \slGapA{}, \slGapB{}, \slGapC{} and \slGapD{} percent at the four tail indices, so the estimator is biased at $\xi=0.6$ even with infinite data, a bias the theory of dispersion does not contain and the measurement shows in the median ratio. Fourth, the limit $T^\ast$ of \cref{cor:history} assumes the same normal approximation and a stationary process.

\begin{conjecture}
\label{conj:bias}
In samples of the size of the experiment the median of $\hat q/q$ is below one for the single loss approximation, by more than the asymptotic bias of the approximation at the true parameters accounts for.
\end{conjecture}

The evidence is the median ratio of the single loss approximation, below one in \slBelowCells{} of the 16 cells, with a mean over cells of \slMeanMed{}; the finite-sample behaviour of the maximum likelihood estimates of the generalised Pareto parameters was compared with other estimators by \citet{hosking1987}, and a bias of the fitted index translates into a bias of the capital through the sensitivity $g_\xi$.
We do not prove the conjecture; \cref{prop:dispersion} gives the spread around the median and not the median.
